\section*{Chapitre \ref{ch:rc:modelling-overnight-rate-products} --- Modélisation des produits à taux overnight}

\begin{solution}{exo:rc:modelling-overnight-rate-products:1}
Un taux interbancaire pour la période en cours était fixé à son début, donc son caplet avait
un gain connu. Un taux overnight composé pour le trimestre en cours n'est connu qu'à
la fin du trimestre~: les fixings restant à venir font du caplet une véritable option.
\end{solution}

\begin{solution}{exo:rc:modelling-overnight-rate-products:2}
Par $\sigma^2\delta/3 = \sigma^2\times0.0833$~: variance totale $\sigma^2\times2.0833$ contre
$\sigma^2\times2$, soit 4.2~\% de plus.
\end{solution}

\begin{solution}{exo:rc:modelling-overnight-rate-products:3}
Le 27 octobre 2027 se situe à 1.077~an de l'échéance spot et le 8 décembre à 1.192~an~; le trimestre se termine à
1.25~: $(1.25-1.077)/0.25 = 69~\%$ et $(1.25-1.192)/0.25 = 23~\%$.
\end{solution}

\begin{solution}{exo:rc:modelling-overnight-rate-products:4}
$\frac{\delta^3/3}{\delta^2S+\delta^3/3} = \frac{\delta/3}{S+\delta/3} = \frac{0.0833}{0.3333} = 25~\%$ pour
$S=0.25$, comme dans le texte (25.2~\% avec retour à la moyenne).
\end{solution}

\begin{solution}{exo:rc:modelling-overnight-rate-products:5}
Les jours de réunion~: son incertitude résiduelle baisse par paliers le 27 octobre et le
8 décembre, et après le 8 décembre elle disparaît. Son theta se concentre sur ces jours
(faible entre-temps), contrairement à la décroissance régulière du modèle généralisé~; un trader doit
couvrir et valoriser le produit réunion par réunion.
\end{solution}

\begin{solution}{exo:rc:modelling-overnight-rate-products:6}
Un taux à terme se calcule à partir de dérivés sur le taux overnight~; si les dérivés sur
le taux à terme prenaient une ampleur importante, il reposerait en partie sur les transactions sur lui-même et pourrait être déplacé
par les flux qu'il est censé mesurer. Restreindre ses dérivés aux
couvertures des utilisateurs finaux préserve la liquidité des produits overnight. Un desk qui vend des caps sur taux à terme
se couvre avec des options sur taux overnight et ne conserve que la variance intra-période et
l'écart entre le taux à terme publié et le taux composé réalisé.
\end{solution}

\begin{solution}{exo:rc:modelling-overnight-rate-products:7}
Le premier caplet, 13\,623~USD, n'existe que dans le cap composé~; les sept autres
caplets valent 24\,599~USD de plus au total que leurs jumeaux prospectifs, du fait de
leur variance intra-période~: 13\,623 plus 24\,599 font 38\,222.
\end{solution}

\begin{solution}{exo:rc:modelling-overnight-rate-products:8}
Les deux hypothèses sont fausses pour les taux composés~: le caplet de la période en cours
reste une option et doit être valorisé, et la variance de chaque caplet court jusqu'à la fin de
sa période, et non à son début. Le pricer sous-évalue chaque cap, ici de 38\,222~USD
sur 100~millions de dollars pendant deux ans, et son vega porte sur les mauvaises dates.
\end{solution}

\begin{solution}{pb:rc:modelling-overnight-rate-products:1}
\textbf{1.} 288\,963~USD sur le SOFR composé, 250\,741~USD sur le taux à terme.
\textbf{2.} Il inclut le caplet du trimestre en cours, et chaque caplet porte un risque
jusqu'à la fin de sa période.
\textbf{3.} 13\,623~USD.
\textbf{4.} Le cap sur SOFR composé, puisque le prêt paie le SOFR composé~; le cap sur taux à
terme laisse subsister l'écart entre le taux du prêt et le taux à terme.
\textbf{5.} 25.2~\% et 4.8~\%.
\textbf{6.} 40\,944~USD par la formule fermée~; la simulation concorde à un cinquième
d'erreur type près.
\textbf{7.} La grille temporelle a manqué les dates de la période (312.5 pas tronqués à 312), si bien que
la période simulée était plus courte que le trimestre et le taux trop bas~: un biais de 6~\% avec une
erreur type serrée.
\textbf{8.} Une variance plus grande de $\delta/(3S) = 8.3~\%$~; Hull--White donne une part intra-période
de 7.9~\%, soit 8.6~\% de plus, presque la même chose.
\textbf{9.} La variance intra-période de Hull--White est proche de $\sigma^2\delta^3/3$, exactement ce que
produit une décroissance linéaire de la volatilité du taux.
\textbf{10.} Le niveau de volatilité et, pour ce strike hors de la monnaie, les
taux à terme.
\textbf{11.} 27 octobre et 8 décembre 2027.
\textbf{12.} Après la décision du 8 décembre, trois semaines avant la fin du trimestre.
\textbf{13.} Zéro~: il ne reste aucune incertitude.
\textbf{14.} Avec des contrats à terme sur les mois de réunion et des options sur ceux-ci, dimensionnés au
poids de chaque réunion dans le trimestre.
\textbf{15.} Le risque~: quand le theta est acquis et où se situe le vega.
\textbf{16.} Le cap sur SOFR composé, qui correspond au prêt~; les 38\,222~USD supplémentaires achètent
une protection sur le trimestre en cours et suppriment le risque de base.
\textbf{17.} L'écart entre le SOFR à terme et le SOFR composé sur chaque trimestre, y compris la
variance intra-période que les options sur taux à terme ne peuvent pas couvrir.
\textbf{18.} Les deux limitent les intérêts du prêt~; le cap sur taux à terme laisse une petite base
qu'ils tolèrent.
\textbf{19.} Résultat retenu~: \emph{le cap SOFR}~: 288\,963~USD sur le SOFR composé contre
250\,741~USD sur le taux à terme~; l'écart de 38\,222~USD correspond au caplet du trimestre en
cours (13\,623) et à la variance intra-période (24\,599).
\textbf{20.} Un risque qui continue de s'accumuler jusqu'à la fin de la période.
\end{solution}

\begin{solution}{iq:rc:modelling-overnight-rate-products:1}
L'Euribor se fixe au début de la période, donc le risque de son caplet s'éteint à ce moment~;
le €STR composé se fixe à la fin, si bien que son caplet porte un risque pendant toute la
période de calcul, vaut davantage, et le caplet de la période en cours est lui aussi une option.
\iqlookfor{le calendrier de fixing et sa conséquence sur la variance.}
\end{solution}

\begin{solution}{iq:rc:modelling-overnight-rate-products:2}
Modéliser le forward du \omterm{def:rc:modelling-overnight-rate-products:rates}{taux rétrospectif} pour chaque période jusqu'à la
fin de celle-ci, avec une volatilité qui décroît vers zéro pendant la période (le \omterm{def:rc:modelling-overnight-rate-products:gfmm}{modèle
de marché forward généralisé})~; avant la période, il coïncide avec le forward prospectif, si bien que
un seul processus donne les deux types de taux.
\iqlookfor{forwards vivants jusqu'à la fin de la période, et volatilité décroissante.}
\end{solution}

\begin{solution}{iq:rc:modelling-overnight-rate-products:3}
Écrire le gain sous la forme $(e^{-\int_0^Sr}-c\,e^{-\int_0^Er})^+$, changer de mesure avec
densité $e^{-\int_0^Sr}/P(0,S)$, et valoriser un put de Black sur $e^{-\int_S^Er}$, lognormal de
moyenne $P(0,E)/P(0,S)$. Sa variance est la variance du prix du zéro-coupon en $S$ plus la
variance de l'intégrale du taux sur la période, $\approx\sigma^2\delta^3/3$.
\iqlookfor{le changement de mesure et les deux termes de variance.}
\end{solution}

\begin{solution}{iq:rc:modelling-overnight-rate-products:4}
Le taux overnight ne bouge que les jours de réunion, si bien que l'incertitude d'une option se
résout par paliers~: le vega et le theta se concentrent sur les réunions, et après la dernière
réunion de la période, le taux de celle-ci est connu. La couverture se fait réunion par réunion, avec des contrats à terme et leurs options.
\iqlookfor{risque par paliers et couverture réunion par réunion.}
\end{solution}

\begin{solution}{iq:rc:modelling-overnight-rate-products:5}
Discrétisation avant statistique~: la grille de simulation contient-elle le début et la
fin de la période (un nombre de pas tronqué raccourcit la période), la convention de
capitalisation, puis la convergence en fonction du pas, les forwards de la courbe utilisés pour la
dérive, et l'actualisation~; comparer avec le cas prospectif, qui admet lui aussi une
forme fermée.
\iqlookfor{biais contre bruit, et la grille.}
\end{solution}

\begin{solution}{iq:rc:modelling-overnight-rate-products:6}
Il se calcule à partir de dérivés sur le taux overnight~; une forte activité de dérivés sur lui
serait circulaire et amincirait le marché sous-jacent. Les recommandations officielles limitent son usage
dérivé aux utilisateurs finaux qui couvrent des prêts indexés sur le taux à terme~; les banques couvrent ceux-ci avec des
instruments sur taux overnight et portent la base.
\iqlookfor{circularité et localisation de la liquidité.}
\end{solution}
